\section{Specialized sources：Wikipedia、GitHub、arXiv}

Wikipedia 是高质量知识源：范围有 notability 和 reliable sources 约束，定期 dumps 可直接下载，无需 crawl。但它不是完美源：编辑集中在少数贡献者，且 data poisoning 可能在 dump 前短暂注入恶意内容。

GitHub 对代码能力重要，也可能帮助 reasoning。其数据包括 repository 内容和 metadata，例如 issues、pull requests、comments。代码数据必须处理 forks、重复、license、malware、PII 和 generated files。

arXiv 提供论文 metadata、PDF 和可选 LaTeX source。它适合科学/数学/CS 内容，但不是 peer review，且许可有 all rights reserved 和 Creative Commons 等差异。

\begin{warningbox}{课堂提示：高质量来源也可能被投毒}
老师用 Wikipedia dump 的时间窗口说明 data poisoning：攻击者可以在定期 dump 前短暂注入恶意编辑，即使管理员随后回滚，污染内容仍可能进入快照。这里的教训不是“Wikipedia 不可信”，而是任何 source 都要记录抓取时间、版本与异常检测；高平均质量并不能消除针对性攻击。
\end{warningbox}

\begin{importantbox}{源不等于数据集}
Wikipedia、GitHub、arXiv 都是 sources。真正进入模型的是经过选择、清洗、格式转换、去重、分片和混合权重设定后的 datasets。数据集设计是模型设计的一部分。
\end{importantbox}

\begin{knowledgebox}{三类专门数据源的互补性}
\begin{tabular}{L{0.18\textwidth}L{0.32\textwidth}L{0.4\textwidth}}
\toprule
来源 & 强项 & 局限 \\
\midrule
Wikipedia & 高密度百科知识、结构清晰、多语 dumps。 & 不是原创思想库，主题受 notability 和编辑群体影响。 \\
GitHub & 代码、工程过程、issues/PR/comments。 & license、fork/duplicate、malware、PII、bot activity。 \\
arXiv & 科研论文、LaTeX source、metadata permissive。 & 未同行评审，PDF/LaTeX 解析复杂，领域分布不均。 \\
\bottomrule
\end{tabular}
\end{knowledgebox}

